\documentclass[11pt,a4paper]{article}
\usepackage[T1]{fontenc}
\usepackage[utf8]{inputenc}
\usepackage[main=italian,provide=*]{babel}
\usepackage{amsmath,amssymb}
\usepackage{geometry}
\geometry{margin=2.3cm,top=2.5cm,bottom=2.7cm}
\usepackage{booktabs}
\usepackage[table,dvipsnames]{xcolor}
\usepackage{tcolorbox}
\tcbuselibrary{skins,breakable}
\usepackage{titlesec}
\usepackage{enumitem}
\usepackage{needspace}
\usepackage{fancyhdr}
\usepackage{graphicx}
\usepackage{hyperref}
\hypersetup{colorlinks=true,linkcolor=Sepia,citecolor=Sepia,urlcolor=Sepia}
\newcommand{\Tr}{\operatorname{Tr}}
\newcommand{\ket}[1]{\lvert #1 \rangle}
\newcommand{\bra}[1]{\langle #1 \rvert}

\definecolor{inkblue}{HTML}{1B3A5C}
\definecolor{lightblue}{HTML}{EAF1F8}
\definecolor{accent}{HTML}{B0562A}
\definecolor{softgray}{HTML}{6B6B6B}
\definecolor{okgreen}{HTML}{2E6B3E}
\definecolor{okgreenbg}{HTML}{EAF5EC}
\definecolor{warnred}{HTML}{9E3B3B}
\definecolor{warnredbg}{HTML}{FBEDED}

\titleformat{\section}
  {\normalfont\Large\bfseries\color{inkblue}}
  {\thesection}{0.6em}{}[{\color{inkblue!40}\titlerule}]
\titlespacing{\section}{0pt}{0.9em}{0.4em}
\titleformat{\subsection}
  {\normalfont\large\bfseries\color{accent}}
  {}{0em}{}
\titlespacing{\subsection}{0pt}{0.6em}{0.2em}

\pagestyle{fancy}
\fancyhf{}
\renewcommand{\headrulewidth}{0.4pt}
\fancyhead[L]{\small\color{softgray}Correlatori dinamici sotto rumore}
\fancyhead[R]{\small\color{softgray}\thepage}
\fancyfoot[C]{\small\color{softgray}Tesi triennale, Samuele --- Relatore: Prof.\ Alessandro Chiesa}

\newtcolorbox{keyeq}[1][]{
  colback=lightblue, colframe=inkblue!70, boxrule=0.6pt,
  arc=2pt, left=8pt, right=8pt, top=3pt, bottom=3pt, #1
}
\newtcolorbox{okbox}[1]{
  colback=okgreenbg, colframe=okgreen!60, boxrule=0.6pt,
  arc=2pt, left=7pt, right=7pt, top=2pt, bottom=2pt,
  fonttitle=\bfseries\color{okgreen}, title={#1}
}
\newtcolorbox{warnbox}[1]{
  colback=warnredbg, colframe=warnred!60, boxrule=0.6pt,
  arc=2pt, left=7pt, right=7pt, top=2pt, bottom=2pt,
  fonttitle=\bfseries\color{warnred}, title={#1}
}

\title{\vspace{-1.5em}\color{inkblue}Correlatori dinamici sotto rumore\\[0.2em]
\Large\normalfont\color{softgray}Il test di Hadamard, il readout, e il compromesso Trotter/rumore}
\author{Note di lavoro --- Tesi triennale, Samuele\\ Relatore: Prof.\ Alessandro Chiesa}
\date{}

\begin{document}
\sloppy
\maketitle
\thispagestyle{fancy}
\vspace{-2.0em}

File di codice: \texttt{correlatori\_rumorosi\_dimero.py}. Punto di
lavoro: $b/J=0.35$, $D/J=0.80$, $J=1$, $t=2$. Correlatore di
riferimento: $C_{21}^{xx}(t)$.

\tableofcontents

\section{Obiettivo}

Un correlatore dinamico $C_{ij}^{\alpha\beta}(t)=\bra{\psi_0}\sigma_i^\alpha(t)\sigma_j^\beta(0)\ket{\psi_0}$
non è un osservabile diretto: si estrae da un circuito di Hadamard test,
che introduce un qubit ancilla e --- a differenza della fedeltà del passo
precedente --- una \textbf{misura reale}, quindi il readout entra in gioco
per la prima volta in questa pipeline. Questo documento costruisce il
circuito rumoroso, deriva la correzione di readout, e trova il
compromesso Trotter/rumore ($N^*$) per questo osservabile specifico.

\section{Il circuito: schema e teoria}

\begin{keyeq}
\centering
$C_{ij}^{\alpha\beta}(t) = \bra{\psi_0}U(t)^\dagger\, \sigma_i^\alpha\, U(t)\, \sigma_j^\beta\ket{\psi_0}$
\end{keyeq}
Schema del circuito (derivazione completa in
\texttt{circuito\_correlatori\_spiegato.tex}):
\begin{enumerate}[leftmargin=1.4em,itemsep=0.15em]
\item preparazione $\ket{\psi_0}$ sul registro (ampiezze esatte, o ansatz
VQE --- vedi Sez.~\ref{sec:ansatz-params});
\item Hadamard sull'ancilla;
\item controlled-$W$ ($W=\sigma_j^\beta$, controllo pieno) sul sito $j$,
\emph{prima} di $U(t)$;
\item $U(t)=e^{-iHt}$ sul registro, non controllato (Trotter, come nel
passo precedente);
\item anti-controlled-$V$ ($V=\sigma_i^\alpha$, attivo su ancilla$=\ket0$)
sul sito $i$, \emph{dopo} $U(t)$;
\item rotazione di base sull'ancilla ($H$ per la parte reale,
$R_x(\pi/2)$ per la parte immaginaria) e misura $Z$.
\end{enumerate}

\subsection{Preparazione: ampiezze esatte o ansatz VQE reale}
\label{sec:ansatz-params}

A differenza del passo precedente, qui la preparazione può usare le
ampiezze esatte (\texttt{ansatz\_params=None}) oppure il circuito VQE
reale (\texttt{ansatz\_params=$\theta^*$}), mirror del pattern già in uso
per anello e catena. Questa opzione, verificata di nuovo in questa
sessione dopo una regressione temporanea del modulo che la implementa
(\texttt{circuito\_correlazioni\_dimero.py}), è la base del confronto di
coerenza in Sez.~\ref{sec:verifica-coerenza}.

\section{Il readout: correzione simmetrica}

\begin{keyeq}
\centering
$\langle Z\rangle_\text{readout} = (1-2p)\,\langle Z\rangle_\text{ideale}$
\end{keyeq}
(caso simmetrico; il caso asimmetrico è trattato separatamente,
\texttt{teoria\_readout\_asimmetrico.tex}). Un fattore puramente
moltiplicativo: importante più avanti per capire perché $N^*$ non dipenda
da $p_\text{readout}$ (visto nello scan sui parametri di rumore).

\section{Verifica di coerenza con la Parte 1}
\label{sec:verifica-coerenza}

\begin{warnbox}{Una regressione trovata e richiusa}
Il modulo \texttt{circuito\_correlazioni\_dimero.py}, che implementa la
preparazione via ansatz VQE reale per il confronto di coerenza qui sotto,
era regredito a una versione precedente (placeholder con
\texttt{NotImplementedError} al posto dell'implementazione) --- non a
causa di un errore mai corretto, ma di un file caricato per sbaglio sopra
la versione giusta (episodio simile, con la stessa causa, già accaduto e
richiuso in una sessione precedente: vedi \texttt{log\_decisioni.md}).
Ricaricato il file corretto e riverificato da zero: tutti i numeri
dichiarati qui sotto sono stati ricalcolati con il modulo funzionante,
non assunti dalla sessione precedente.
\end{warnbox}

Confronto, a rumore nullo, fra il correlatore calcolato dal circuito
completo (preparazione ansatz VQE + Trotter + Hadamard test) e quello
calcolato dal modulo di riferimento della Parte 1 (stessa preparazione
ansatz, nessun rumore):

\begin{center}
\renewcommand{\arraystretch}{1.25}
\begin{tabular}{@{}lccc@{}}
\toprule
$\boldsymbol{(i,\alpha,j,\beta,t,N)}$ & \textbf{Parte 1} & \textbf{Rumore nullo} & $\boldsymbol{|\Delta|}$\\
\midrule
$(2,x,1,x,2.0,20)$ & $+0.074599+0.257465i$ & $+0.074599+0.257465i$ & $2.5\times10^{-14}$\\
$(1,z,1,z,1.0,20)$ & $-0.158218+0.969111i$ & $-0.158218+0.969111i$ & $4.2\times10^{-15}$\\
$(1,x,1,y,0.5,20)$ & $+0.205804+0.228287i$ & $+0.205804+0.228287i$ & $5.0\times10^{-14}$\\
\bottomrule
\end{tabular}
\end{center}

\begin{okbox}{Verificato}
Coincidenza a precisione quasi macchina --- il circuito rumoroso, a
rumore spento, riproduce esattamente il calcolo di riferimento. Ulteriori
tre controlli specifici sul modulo di preparazione VQE: infedeltà
ansatz/esatto $1.6\times10^{-11}$ (dipende dal seed del VQE, non un
problema); consistenza fra le due preparazioni sui tre correlatori
sopra, $3\times10^{-7}$--$4\times10^{-6}$; costo in gate, $18$ CNOT
(preparazione esatta) contro $19$ CNOT (preparazione VQE) a $N=5$ ---
un solo CNOT in più per il circuito VQE, trascurabile.
\end{okbox}

\section{Conteggio gate: crescita lineare}

\begin{center}
\renewcommand{\arraystretch}{1.2}
\begin{tabular}{@{}cc@{}}
\toprule
$\boldsymbol{N}$ & \textbf{CNOT totali}\\
\midrule
$1$ & $7$\\
$2$ & $10$\\
$5$ & $19$\\
$10$ & $34$\\
$20$ & $64$\\
$40$ & $124$\\
\bottomrule
\end{tabular}
\end{center}

\noindent Incremento costante di 3 CNOT per unità di $N$ (il contributo
di un passo di Trotter) --- nessuna fusione spuria, la transpilazione a
blocchi funziona come deve.

\section{Formula di readout: Monte Carlo contro analitica}

\begin{center}
\renewcommand{\arraystretch}{1.2}
\begin{tabular}{@{}lc@{}}
\toprule
\textbf{Quantità} & \textbf{Valore}\\
\midrule
$\langle Z\rangle$ analitico (rumore di gate + correzione readout) & $+0.044089$\\
$\langle Z\rangle$ Monte Carlo ($2\times10^5$ shot, readout vero) & $+0.046550$\\
$|$differenza$|$ (atteso $\lesssim0.0067$, $3\sigma$) & $2.46\times10^{-3}$\\
\bottomrule
\end{tabular}
\end{center}

\noindent Due cammini di calcolo indipendenti (formula analitica contro
simulazione Monte Carlo con \texttt{ReadoutError} vero) coincidono entro
l'errore statistico atteso.

\section{Risultato: la curva \texorpdfstring{$|C(N)|$}{|C(N)|} e \texorpdfstring{$N^*$}{N*}}

\begin{figure}[h]
\centering
\includegraphics[width=0.72\linewidth]{figure/fig_passo4_correlatore_vs_N.pdf}
\caption{\textit{Modulo del correlatore rumoroso $|C_{21}^{xx}(t{=}2,N)|$
in funzione di $N$. Massimo a $N^*=5$ --- più basso dell'$N^*=8$ della
fedeltà di Trotter (atteso: il circuito di Hadamard test ha più gate a
parità di $N$, per via dei due blocchi controllati aggiuntivi, quindi il
rumore satura prima).}}
\end{figure}

\begin{center}
\renewcommand{\arraystretch}{1.2}
\begin{tabular}{@{}lc@{}}
\toprule
\textbf{Quantità} & \textbf{Valore}\\
\midrule
$N^*$ & $5$\\
$|C(N^*)|$ & $0.250002$\\
\bottomrule
\end{tabular}
\end{center}

\section{Considerazioni e cosa serve al passo successivo}

Tre fatti da qui sono usati esplicitamente nello scan sui parametri di
rumore: (1) $N^*=5$ è specifico di questo correlatore, a questi
parametri --- lo scan ne verifica la stabilità al variare di
$\varepsilon_{1q},\varepsilon_{2q},p_\text{readout}$, non la dà per
scontata; (2) la correzione di readout è un fattore puramente
moltiplicativo, $(1-2p)$: questo è l'argomento analitico (verificato,
non solo dichiarato) per cui $N^*$ non può dipendere da $p_\text{readout}$
per readout simmetrico --- moltiplicare per una costante positiva non
sposta la posizione di un massimo; (3) il conteggio gate lineare in $N$
è la base per stimare come il rumore accumulato cresca al variare dei
parametri del canale.

\begin{okbox}{Sintesi in una riga}
$N^*=5$ per il correlatore dinamico, verificato con il modulo di
preparazione VQE ora funzionante dopo una regressione richiusa;
correzione di readout puramente moltiplicativa, la chiave per capire
perché $N^*$ non dipenderà da $p_\text{readout}$.
\end{okbox}

\needspace{4\baselineskip}
\section*{File prodotti}
\addcontentsline{toc}{section}{File prodotti}

\texttt{correlatori\_rumorosi\_dimero.py}, questo documento. Modulo di
preparazione VQE: \texttt{circuito\_correlazioni\_dimero.py} (ricaricato
e riverificato in questa sessione).

\end{document}
